% (User input window)
% 2026-02-14 15:08:47（Asia/Singapore）

\paragraph{Finite Laplace Principle and Phase Diagram: Leading Layer, Leading Coefficient and Symmetric Group Gibbs Limit}
Following the above setup and notation.
To simplify the exposition, for any partition \(\nu\vdash q\) we define its \textbf{partition pressure} and \textbf{partition leading coefficient}
\[
\mathcal{P}(\nu):=\sum_{r=1}^{q} c_r(\nu)\,P(r),
\qquad
K(\nu):=\prod_{r=1}^{q}\kappa_r^{\,c_r(\nu)}.
\]
We also denote the spectral gap parameter
\[
\varepsilon_*:=\min_{1\le r\le q}\log\!\left(\frac{\rho_r}{\theta_r}\right)>0
\]
(where \(\rho_r=e^{P(r)}\), and \(0<\theta_r<\rho_r\) is given by \eqref{eq:pom-sr-main-term-kappa}).

\begin{lemma}[Exponential Leading Term of Partition Monomial]\label{lem:pom-partition-monomial-exponential-main-term}
For any fixed partition \(\nu\vdash q\) we have
\[
\mathsf{M}_\nu(m)=K(\nu)\,e^{m\mathcal{P}(\nu)}+O\!\left(e^{m(\mathcal{P}(\nu)-\varepsilon_*)}\right)\qquad(m\to\infty).
\]
\end{lemma}
\begin{proof}
By \eqref{eq:pom-sr-main-term-kappa},
\(
S_r(m)=\kappa_r\rho_r^{\,m}+O(\theta_r^{\,m})
=\kappa_r e^{mP(r)}\bigl(1+O(e^{-m\varepsilon_*})\bigr)
\)
holds simultaneously for all \(1\le r\le q\).
Substituting into \(\mathsf{M}_\nu(m)=\prod_{r=1}^{q}S_r(m)^{c_r(\nu)}\) and taking the finite product, the error can be absorbed into constants and the relative error remains \(O(e^{-m\varepsilon_*})\).
\end{proof}

\begin{theorem}[Finite Laplace Principle for Schur Channel Trace: Decidable Formula for Leading Exponent and Leading Coefficient]\label{thm:pom-schur-trace-finite-laplace-principle}
Fix \(q\ge 2\), and follow the partition summation formula of Proposition \ref{prop:pom-schur-tomography-by-partitions}
\[
T_{q,\lambda}(m)=f^\lambda\sum_{\nu\vdash q}\frac{\chi^\lambda(\nu)}{z_\nu}\,\mathsf{M}_\nu(m).
\]
For any \(\lambda\vdash q\) and any exponential layer \(v\in\RR\) define the layer coefficient
\[
A_\lambda(v):=
f^\lambda\!\!\sum_{\substack{\nu\vdash q\\ \mathcal{P}(\nu)=v}}\frac{\chi^\lambda(\nu)}{z_\nu}\,K(\nu),
\qquad
\mathcal{V}_\lambda:=\left\{\mathcal{P}(\nu):\nu\in\mathrm{Supp}(\lambda)\right\}.
\]
If there exists \(v\in\mathcal{V}_\lambda\) such that \(A_\lambda(v)\neq 0\), let
\[
\mathcal{P}_\lambda^{\mathrm{eff}}
:=
\max\{v\in\mathcal{V}_\lambda:\ A_\lambda(v)\neq 0\}.
\]
Then there exists \(\varepsilon_\lambda>0\) such that
\[
T_{q,\lambda}(m)
=
A_\lambda(\mathcal{P}_\lambda^{\mathrm{eff}})\,e^{m\mathcal{P}_\lambda^{\mathrm{eff}}}
+O\!\left(e^{m(\mathcal{P}_\lambda^{\mathrm{eff}}-\varepsilon_\lambda)}\right),
\qquad
\limsup_{m\to\infty}\frac{1}{m}\log|T_{q,\lambda}(m)|=\mathcal{P}_\lambda^{\mathrm{eff}}.
\]
If for all \(v\in\mathcal{V}_\lambda\) we have \(A_\lambda(v)=0\), then
\[
\limsup_{m\to\infty}\frac{1}{m}\log|T_{q,\lambda}(m)|
\le
\max_{v\in\mathcal{V}_\lambda}v-\varepsilon_*,
\]
i.e., the exponential leading order drops by at least \(\varepsilon_*\).
In particular, if \(\mathcal{P}_\lambda\) and \(\mathcal{M}_\lambda\) are as defined above and \(A_\lambda(\mathcal{P}_\lambda)\neq 0\), then \(\mathcal{P}_\lambda^{\mathrm{eff}}=\mathcal{P}_\lambda\), consistent with the non-degenerate case of Theorem \ref{thm:pom-schur-channel-radius-character-sieve}.
\end{theorem}
\begin{proof}
By Lemma \ref{lem:pom-partition-monomial-exponential-main-term}, we can write \(T_{q,\lambda}(m)\) as a finite sum of exponential terms, and re-stratify by \(\mathcal{P}(\nu)\):
\[
T_{q,\lambda}(m)=\sum_{v\in\mathcal{V}_\lambda} A_\lambda(v)\,e^{mv}+O\!\left(e^{m(\max\mathcal{V}_\lambda-\varepsilon_*)}\right).
\]
If there exists a nonzero layer coefficient, taking \(v=\mathcal{P}_\lambda^{\mathrm{eff}}\) as the maximal nonzero layer yields the leading term and exponent rate;
the error term is controlled jointly by the exponential gap from lower layers \(v<\mathcal{P}_\lambda^{\mathrm{eff}}\) and the uniform deficit in the lemma.
If all layer coefficients are zero, then the leading-term layer is completely canceled, leaving only error contributions from each \(\mathsf{M}_\nu(m)\), yielding the stated upper bound.
\end{proof}

\paragraph{Total Channel Cancellation of Resonant Layer: Character Orthogonality and Rigidity}
Denote the global maximum pressure and its maximizing set by
\[
\mathcal{P}_{\max}:=\max_{\nu\vdash q}\mathcal{P}(\nu),
\qquad
\mathcal{R}_{\max}:=\{\nu\vdash q:\ \mathcal{P}(\nu)=\mathcal{P}_{\max}\}.
\]
Introduce the resonant layer weight function (viewed as a class function on \(S_q\))
\[
w(\nu):=K(\nu)\,\mathbf 1_{\{\nu\in\mathcal{R}_{\max}\}}.
\]

\begin{proposition}[Resonant Layer Fourier Coefficient and Necessary and Sufficient Condition for Total Channel Cancellation]\label{prop:pom-schur-resonant-layer-total-cancellation}
For each \(\lambda\vdash q\) define the resonant layer Fourier coefficient
\[
b_\lambda:=\sum_{\nu\vdash q}\frac{\chi^\lambda(\nu)}{z_\nu}\,w(\nu)
=\sum_{\nu\in\mathcal{R}_{\max}}\frac{\chi^\lambda(\nu)}{z_\nu}\,K(\nu).
\]
Then there holds the inversion formula: for any \(\mu\vdash q\),
\[
\boxed{\;
w(\mu)=\sum_{\lambda\vdash q}\chi^\lambda(\mu)\,b_\lambda\; }.
\]
Thus the following two propositions are equivalent:
\begin{enumerate}
  \item For all non-trivial \(\lambda\neq(q)\) we have \(b_\lambda=0\);
  \item \(w(\nu)\) is constant on all conjugacy classes (i.e., \(w\) is orthogonal to all non-trivial irreducible characters).
\end{enumerate}
In particular, if \(\mathcal{R}_{\max}\subsetneq\{\nu\vdash q\}\), then there must exist some non-trivial \(\lambda\neq(q)\) such that \(b_\lambda\neq 0\), i.e., the resonant exponential layer cannot be simultaneously canceled by all non-trivial Schur channels.
\end{proposition}
\begin{proof}
By the column orthogonality relation for symmetric group characters
\(
\sum_{\lambda\vdash q}\chi^\lambda(\mu)\chi^\lambda(\nu)=z_\mu\,\mathbf 1_{\{\mu=\nu\}}
\),
we have
\[
\sum_{\lambda\vdash q}\chi^\lambda(\mu)\,b_\lambda
=\sum_{\nu\vdash q}\frac{w(\nu)}{z_\nu}\sum_{\lambda\vdash q}\chi^\lambda(\mu)\chi^\lambda(\nu)
=\sum_{\nu\vdash q}\frac{w(\nu)}{z_\nu}\,z_\mu\,\mathbf 1_{\{\mu=\nu\}}
=w(\mu),
\]
so the inversion formula holds.
If for all \(\lambda\neq(q)\) we have \(b_\lambda=0\), then \(w(\mu)=\chi^{(q)}(\mu)\,b_{(q)}\equiv b_{(q)}\) is constant;
conversely, if \(w\) is a constant class function, then it is orthogonal to all non-trivial characters, so \(b_\lambda=0\) (\(\lambda\neq(q)\)).
Finally, if \(\mathcal{R}_{\max}\) is not the set of all partitions, then \(w\) is not a constant class function, hence at least one non-trivial Fourier coefficient must be nonzero.
\end{proof}

\begin{theorem}[Pressure Vector Fan Phase Diagram: Piecewise Linear Structure of Channel Leading Layer]\label{thm:pom-schur-pressure-fan-diagram}
Fix \(q\ge2\), and denote the pressure vector by \(p=(P(1),\dots,P(q))\in\RR^q\).
For each partition \(\nu\vdash q\), \(\mathcal{P}(\nu)=\langle c(\nu),p\rangle\) is a linear function of \(p\) (where \(c(\nu)=(c_1(\nu),\dots,c_q(\nu))\)).
Let the hyperplane family
\[
\mathcal{H}_q
:=
\left\{\langle c(\nu)-c(\mu),p\rangle=0:\ \nu,\mu\vdash q,\ \nu\neq\mu\right\}.
\]
Then for any \(p\notin\bigcup\mathcal{H}_q\), all \(\{\mathcal{P}(\nu)\}_{\nu\vdash q}\) are pairwise distinct.
Thus for each \(\lambda\vdash q\), the maximizing set \(\mathcal{M}_\lambda\) is a singleton \(\{\nu_\lambda(p)\}\), and there exists \(\varepsilon_\lambda(p)>0\) such that
\[
T_{q,\lambda}(m)
=
f^\lambda\,\frac{\chi^\lambda(\nu_\lambda(p))}{z_{\nu_\lambda(p)}}\,K(\nu_\lambda(p))\,
e^{m\mathcal{P}(\nu_\lambda(p))}
+O\!\left(e^{m(\mathcal{P}(\nu_\lambda(p))-\varepsilon_\lambda(p))}\right),
\]
where we can take
\[
\varepsilon_\lambda(p)=\min\!\left\{
\varepsilon_*,
\min_{\nu\in\mathrm{Supp}(\lambda),\,\nu\neq\nu_\lambda(p)}\bigl(\mathcal{P}(\nu_\lambda(p))-\mathcal{P}(\nu)\bigr)
\right\}>0.
\]
\end{theorem}
\begin{proof}
When \(p\notin\bigcup\mathcal{H}_q\), \(\mathcal{P}(\nu)\) is strictly totally ordered on the finite set, so the maximum for each \(\lambda\) is uniquely attained on \(\mathrm{Supp}(\lambda)\) at \(\nu_\lambda(p)\).
Moreover, \(\nu_\lambda(p)\in\mathrm{Supp}(\lambda)\) yields \(\chi^\lambda(\nu_\lambda(p))\neq 0\), and \(K(\nu_\lambda(p))>0\), so the leading layer coefficient is a single nonzero term.
Substituting into Theorem \ref{thm:pom-schur-trace-finite-laplace-principle} and simultaneously incorporating the uniform deficit from Lemma \ref{lem:pom-partition-monomial-exponential-main-term} yields the stated error exponent.
\end{proof}

\begin{proposition}[Leading Cycle-Length Selection: Knapsack Structure of Maximum Exponential Layer]\label{prop:pom-partition-pressure-knapsack}
For \(1\le r\le q\) define the per-part pressure \(a_r:=P(r)/r\), and denote \(a_*:=\max_{1\le r\le q}a_r\), \(R_*:=\{r:\ a_r=a_*\}\).
Then for any partition \(\nu\vdash q\) we always have
\[
\mathcal{P}(\nu)\le q\,a_*.
\]
If \(\nu\) attains the maximum value \(\mathcal{P}_{\max}\), then any part length \(r\) appearing in \(\nu\) must satisfy \(r\in R_*\).
If \(R_*=\{r_*\}\) is a singleton and \(r_*\mid q\), then the maximizing partition is unique and equals \(\nu=(r_*)^{q/r_*}\).
\end{proposition}
\begin{proof}
By the constraint \(\sum_{r=1}^{q} r\,c_r(\nu)=q\), we have
\[
\mathcal{P}(\nu)=\sum_{r=1}^{q} r\,c_r(\nu)\,a_r\le a_*\sum_{r=1}^{q} r\,c_r(\nu)=q a_*.
\]
If some \(r\notin R_*\) still appears, then the corresponding term is strictly less than \(a_* r c_r(\nu)\), so the total is strictly less than \(q a_*\), cannot attain the maximum.
Finally, when \(R_*=\{r_*\}\) and \(r_*\mid q\), the constraint forces all part lengths to be \(r_*\), and from \(\sum r_*c_{r_*}=q\) we obtain uniqueness.
\end{proof}

\begin{corollary}[Total Channel Collapse Under Identity Sector Dominance]\label{cor:pom-identity-sector-collapse-all-schur-channels}
If the strict dominance holds:
\[
\frac{P(1)}{1}>\max_{2\le r\le q}\frac{P(r)}{r},
\]
then \(\mathcal{P}_{\max}=qP(1)\) and the unique maximizing partition is \(\nu_0=(1^q)\).
Thus for all \(\lambda\vdash q\) there exists \(\varepsilon>0\) such that
\[
T_{q,\lambda}(m)
=
\frac{(f^\lambda)^2}{q!}\,\kappa_1^{\,q}\,e^{m qP(1)}
+O\!\left(e^{m(qP(1)-\varepsilon)}\right).
\]
\end{corollary}
\begin{proof}
By Proposition \ref{prop:pom-partition-pressure-knapsack}, strict dominance means \(R_*=\{1\}\), so the maximizing partition is uniquely \((1^q)\).
Moreover, \(\chi^\lambda(1^q)=f^\lambda>0\), \(z_{1^q}=q!\), \(K(1^q)=\kappa_1^q\), so the leading coefficient is strictly nonzero.
Substituting \(\nu_0\) into Theorem \ref{thm:pom-schur-trace-finite-laplace-principle} yields the stated leading term and exponential error.
\end{proof}

\paragraph{Symmetric Group Gibbs Limit: Conjugacy Class Freezing and Order Parameter}
Define the Gibbs distribution on \(S_q\) (for fixed \(m\)):
\[
\mathbb P_{m,q}(\sigma):=\frac{\mathcal{C}_\sigma(m)}{\sum_{\pi\in S_q}\mathcal{C}_\pi(m)}
=\frac{\mathsf{M}_{\nu(\sigma)}(m)}{\sum_{\pi\in S_q}\mathsf{M}_{\nu(\pi)}(m)}.
\]
Its projection onto the partition space is
\[
\mathbb P^{\mathrm{part}}_{m,q}(\nu)
:=
\sum_{\sigma:\,\nu(\sigma)=\nu}\mathbb P_{m,q}(\sigma)
=
\frac{\frac{q!}{z_\nu}\mathsf{M}_\nu(m)}{\sum_{\mu\vdash q}\frac{q!}{z_\mu}\mathsf{M}_\mu(m)}.
\]
Denote the total variation distance \(D_{\mathrm{TV}}(\mu,\nu):=\frac12\sum_x|\mu(x)-\nu(x)|\) (finite state space).

\begin{theorem}[Freezing and Channel Order Parameter Limit Under Unique Leading Cycle Type]\label{thm:pom-symmetric-group-gibbs-freezing}
If \(\mathcal{R}_{\max}=\{\nu_*\}\) is a singleton, and denote
\[
\delta:=\mathcal{P}_{\max}-\max_{\nu\neq\nu_*}\mathcal{P}(\nu)>0,
\qquad
\delta_*:=\min(\delta,\varepsilon_*)>0,
\]
then there exists a constant \(C_q>0\) such that
\[
D_{\mathrm{TV}}\!\left(\mathbb P_{m,q},\ \mathrm{Unif}(\mathrm{Cl}(\nu_*))\right)\le C_q\,e^{-m\delta_*},
\qquad
D_{\mathrm{TV}}\!\left(\mathbb P_{m,q}^{\mathrm{part}},\ \delta_{\nu_*}\right)\le C_q\,e^{-m\delta_*}.
\]
For any \(\lambda\vdash q\) define the order parameter
\[
\mathfrak m_\lambda(m):=\EE_{m,q}\bigl[\chi^\lambda(\sigma)\bigr],
\]
then we have
\[
\mathfrak m_\lambda(m)=\chi^\lambda(\nu_*)+O(e^{-m\delta_*}),
\qquad
\frac{T_{q,\lambda}(m)}{T_{q,(q)}(m)}=f^\lambda\,\mathfrak m_\lambda(m)
=f^\lambda\,\chi^\lambda(\nu_*)+O(e^{-m\delta_*}).
\]
\end{theorem}
\begin{proof}
By Lemma \ref{lem:pom-partition-monomial-exponential-main-term},
\(
\mathsf{M}_\nu(m)=K(\nu)e^{m\mathcal{P}(\nu)}\bigl(1+O(e^{-m\varepsilon_*})\bigr)
\).
Substituting into the finite ratio formula for \(\mathbb P_{m,q}^{\mathrm{part}}(\nu)\), and controlling the remaining terms by the unique maximum layer gap \(\delta\), we obtain
\(
\mathbb P_{m,q}^{\mathrm{part}}(\nu_*)=1+O(e^{-m\delta_*})
\)
and \(\mathbb P_{m,q}^{\mathrm{part}}(\nu)=O(e^{-m\delta_*})\) (\(\nu\neq\nu_*\)).
Thus the TV distance in partition space exponentially converges to \(\delta_{\nu_*}\).
\par
Since \(\mathcal{C}_\sigma(m)\) is constant within conjugacy classes, conditioned on \(\nu_*\), \(\mathbb P_{m,q}\) is uniform on \(\mathrm{Cl}(\nu_*)\),
hence
\(
D_{\mathrm{TV}}(\mathbb P_{m,q},\mathrm{Unif}(\mathrm{Cl}(\nu_*)))\le 1-\mathbb P_{m,q}^{\mathrm{part}}(\nu_*)
\),
yielding the TV bound on \(S_q\).
\par
Characters are class functions, so \(\mathfrak m_\lambda(m)=\sum_{\nu}\mathbb P_{m,q}^{\mathrm{part}}(\nu)\chi^\lambda(\nu)\), and the convergence of the partition distribution yields the stated limit and error.
Finally, from
\(
T_{q,\lambda}(m)=\frac{f^\lambda}{q!}\sum_{\sigma\in S_q}\chi^\lambda(\sigma)\mathcal{C}_\sigma(m)
\)
and
\(
T_{q,(q)}(m)=\frac{1}{q!}\sum_{\sigma\in S_q}\mathcal{C}_\sigma(m)
\)
we obtain \(T_{q,\lambda}(m)=f^\lambda T_{q,(q)}(m)\mathfrak m_\lambda(m)\).
\end{proof}

\begin{theorem}[Mixture Limit in Multi-Maximizer Phase: Conjugacy Class Mixture and Order Parameter Jump]\label{thm:pom-symmetric-group-gibbs-mixture}
If \(|\mathcal{R}_{\max}|\ge 2\), define the maximum layer weight
\[
W(\nu):=\frac{K(\nu)}{z_\nu}\qquad(\nu\in\mathcal{R}_{\max}),
\qquad
\pi_\infty(\nu):=\frac{W(\nu)}{\sum_{\mu\in\mathcal{R}_{\max}}W(\mu)}.
\]
Then there holds the limit
\[
\lim_{m\to\infty}\mathbb P^{\mathrm{part}}_{m,q}(\nu)=
\begin{cases}
\pi_\infty(\nu),& \nu\in\mathcal{R}_{\max},\\
0,& \nu\notin\mathcal{R}_{\max},
\end{cases}
\]
and
\[
\mathbb P_{m,q}\xrightarrow[m\to\infty]{}\sum_{\nu\in\mathcal{R}_{\max}}\pi_\infty(\nu)\,\mathrm{Unif}(\mathrm{Cl}(\nu)),
\qquad
\lim_{m\to\infty}\mathfrak m_\lambda(m)=\sum_{\nu\in\mathcal{R}_{\max}}\pi_\infty(\nu)\,\chi^\lambda(\nu).
\]
When the pressure vector \(p\) crosses a wall of \(\mathcal{H}_q\) and causes \(\mathcal{R}_{\max}\) to change, the above limiting weights and order parameters exhibit piecewise constant jumps.
\end{theorem}
\begin{proof}
Still by Lemma \ref{lem:pom-partition-monomial-exponential-main-term}, substituting into \(\mathbb P_{m,q}^{\mathrm{part}}(\nu)\) with
\(
\mathsf{M}_\nu(m)=K(\nu)e^{m\mathcal{P}(\nu)}(1+o(1))
\)
and extracting the common factor \(\mathcal{P}_{\max}\):
Non-maximum layers vanish exponentially, maximum layers remain and are normalized by \(W(\nu)\), yielding the partition space limit.
Since weights within classes are constant, the limit on \(S_q\) is a convex combination of uniform distributions over each maximum conjugacy class; the character expectation converges accordingly, yielding the order parameter limit formula.
\end{proof}

